\section{Motivation}

\subsection{Existing System}
\noindent There are several automated systems used for managing final admission processes in educational institutions. These systems differ in features, capabilities, and implementation depending on institutional needs and available resources.

\vspace{0.3cm}

\noindent Currently, JUST provides a web-based application titled ``JUST Admission 2025''. While this system includes some useful functionalities, it is not a fully integrated or complete solution. Several important features are missing, such as admission confirmation, sending admission confirmation messages, and verification of required documents.

\vspace{0.3cm}

\noindent Another issue in the current system is that any user can register regardless of whether they have been selected for admission or not. This creates inconsistency and reduces system reliability. Additionally, the system lacks proper validation mechanisms, which makes it difficult to ensure the authenticity and correctness of submitted data.

\vspace{0.3cm}

\subsection{Limitations of the System}
\noindent The current system suffers from several limitations that affect its efficiency and usability.

\vspace{0.3cm}

\noindent First, there is no proper validation process implemented. As a result, users can input incorrect or incomplete data without restriction. Second, the system allows unrestricted data insertion but does not provide proper mechanisms for updating or removing data, which can lead to data inconsistency.

\vspace{0.3cm}

\noindent Third, the system architecture and code structure are not well organized, making it difficult to maintain, scale, or extend the system in the future.

\vspace{0.3cm}

\noindent Another significant limitation is that users who are not eligible (i.e., those who did not get a chance for admission) can still register in the system. This reduces the accuracy and trustworthiness of the admission process.

\vspace{0.3cm}

\noindent Finally, key features such as admission confirmation, automated message sending, and document verification are either missing or not fully implemented, which prevents the system from functioning as a complete automated admission solution.